% Changed: restore service-object argument and verified serving evidence.
% Protected: qualification semantics, measurement clocks, adverse results and citation keys.
% Meaning-Check: PRESERVED
% Risk-Flags: none
\begin{abstract}
Repeated context reduces LLM serving cost only when a serving engine can reuse its physical KV state. Dynamic questions and observations often interrupt that prefix, while reordering can change generation and expose unsupported state to later tenants. We present ContractKV, a service control plane that governs prompt layout, behavioral evidence, and shared-state publication together. Its compiler constructs dependency-valid stable-prefix layouts; scoped certification checks model-generated token paths and termination; Prefix-Selective Transactional Promotion publishes only the exact, block-aligned certified prefix; and Service-Template Certificates bind authorization, version and lifecycle so admission costs can be shared across tenants. On six source-disjoint HotpotQA and RepoBench templates, Qwen3-8B, Mistral-7B and Llama-3.1-8B achieve 1.619$\times$, 1.109$\times$ and 1.070$\times$ full-cost throughput in three paired fresh-process repetitions under greedy decoding. These zero-slack measurements charge successful and failed qualification and waiting. All 1,296 admitted compiled future pairs retain their raw output token IDs. Separately, deployment-qualified controls reduce verifier calls from 125 to five and unhidden verification cost from 25.05 to 3.24 ms/request. At a fixed 80\%-load point with 25 tenants per template, an open-arrival study raises 500-ms SLO goodput from 0.173 to 0.831 requests/s. These distinct protocols connect behavioral admission to actual service gains: a reusable prefix becomes a governed object whose qualification, visibility and amortization can be managed explicitly.
\end{abstract}
